% \subsection{Shared utilities}
% \label{sec:shared}

% Before walking through ingestion or training, it is worth describing the small shared library I extracted once several packages started duplicating the same path logic. Early prototype version had several hard-coded directories after repetitive copy-paste, I centralised the behaviour under \texttt{libs/shared} making the code reusable across different components.

% Each component has configuration loading as shown in figure~\ref{fig:example_libs}. The main use of configuration loading is to  (\codecite{code_config_loader}) walk up from the package until it finds \texttt{pyproject.toml}, then deep-merge \texttt{config/default.yaml} which a config file with a package-specific file. As a developer, This pattern eased up profoundly in configuration management and let me modify MinIO endpoints or MariaDB paths in one place during Docker bring-up without touching Python imports. None of these modules are glamorous, but they made the pipeline development more efficient and less error-prone.


% Object storage helpers in \codecite{code_s3} uses boto3 against MinIO bucket names, upload helpers, and existence checks. \codecite{code_storage_config} turns a month string like \texttt{2024-05} into the right S3 URI or local MariaDB file path so the CLI and Airflow tasks agree on inputs. Logging (\codecite{code_logger}) appends to \texttt{logs/lichess.log} with rotation; when a DAG task failed at 03:00 I could trace it without digging through container stdout alone. I also added a thin \texttt{LichessException} wrapper (see \repo{libs/shared/exception.py}) so CLI failures print file context with helpful information. Exceptions are not user-facing API errors yet, they are just made for developer friendly at this point.


